\begin{frame}
\titlepage
\end{frame}

\begin{frame}{Wat u na deze inleiding kunt uitleggen}
\begin{itemize}
\item Wat u met MISP deelt en waarom context belangrijk is.
\item Hoe u ontvangers en een bruikbare vervolgstap kiest.
\item Hoe feedback gedeelde informatie verbetert.
\end{itemize}
\note{Een korte inleiding voor analisten, incidentresponders, leidinggevenden en nieuwe deelnemers aan een deelgemeenschap. Reken op 15--20 minuten en vijf minuten discussie. Een account, API-sleutel of technische voorbereiding is niet nodig. Gebruik de fictieve casus door de hele inleiding. De drie taalversies hebben dezelfde opbouw en delen de bewerkbare diagrammen. Ga voor de praktijk verder met de cursus MISP-gebruik.}
\end{frame}

\begin{frame}{Begin met een gezamenlijke vraag}
\begin{itemize}
\item Team A onderzoekt een verdacht aanmeldbericht.
\item Het registreert \IntroCase{} en \texttt{192.0.2.25}.
\item Team B zag hetzelfde domein in een ander rapport.
\item Wat leren ze samen en wat blijft onzeker?
\end{itemize}
\note{Dit is een fictieve casus. Het domein .example en het adresbereik 192.0.2.0/24 zijn gereserveerd voor documentatie; de waarden zijn geen kwaadaardige indicatoren. Teams kunnen werk hergebruiken, waarnemingen vergelijken en context vragen. Een gedeeld domein bewijst geen gezamenlijke campagne of actor. Vraag welke bron, tijdinformatie en betrouwbaarheid nodig zijn voordat iemand handelt.}
\end{frame}

\begin{frame}{MISP maakt waarnemingen herbruikbaar}
\begin{itemize}
\item Open source software om gestructureerde informatie te delen.
\item Een plek om dreigingsinformatie te verzamelen en beoordelen.
\item Samenwerking tussen mensen, tools en teams.
\item Context voor onderzoek, detectie en respons.
\end{itemize}
\note{MISP ondersteunt technische indicatoren en bredere context over incidenten, fraude en dreigingen. Analisten, responders, onderzoekers en deelgemeenschappen hebben verschillende behoeften. Het platform helpt gestructureerd uitwisselen zonder dat iedereen dezelfde interne tools gebruikt. Software ondersteunt het proces, maar schept geen vertrouwen en beoordeelt niet automatisch iedere bron of vervolgstap.}
\end{frame}

\begin{frame}{Bewaar de waarneming samen met de context}
\begin{itemize}
\item Event: de context van het rapport of onderzoek.
\item Attribuut: één getypeerde waarneming, zoals een domein.
\item Object: samenhangende attributen volgens een sjabloon.
\item Tags en Galaxies: classificatie en gestructureerde context.
\end{itemize}
\note{Het event beschrijft in de casus het verdachte aanmeldbericht. Domein- en IP-attributen zijn waarnemingen; een domain-ip-object kan een betekenisvol verband groeperen. Taxonomieën leveren classificatietermen en Galaxy-clusters herbruikbare kennis. Een waarneming is niet vanzelf een inzetklare detectie-indicator. Bewaar herkomst en onzekerheid. De cursussen over gebruik, objecten en Galaxies werken deze begrippen verder uit.}
\end{frame}

\begin{frame}{Van waarnemingen naar gedeelde context}
\IntroDiagram{sharing-flow}
\note{Volg de volle pijlen van bronnen via structureren, beoordelen en deelinstellingen naar gebruik. De stippelpijl brengt waarnemingen, correcties en nieuwe context terug naar de beoordeling. Dit is een conceptuele werkwijze, geen vaste automatische uitvoeringsketen. Feeds en verrijkingsdiensten worden geselecteerd en ingesteld; niet iedere invoer is betrouwbaar en niet iedere uitvoer wordt gepubliceerd. Distributie, rollen en de gekozen uitwisselingsmethode begrenzen wat de instantie verlaat.}
\end{frame}

\begin{frame}{Een correlatie is een aanknopingspunt}
\begin{itemize}
\item Een gedeelde waarde kan relevante rapporten verbinden.
\item Controleer bron, tijdstip en omliggende waarnemingen.
\item Denk aan gedeelde infrastructuur, hergebruik en foutpositieven.
\item Een overeenkomst bewijst geen oorzaak of attributie.
\end{itemize}
\note{Het gezamenlijke domein kan relevant, onschuldig, hergebruikt of verkeerd geregistreerd zijn. Correlatie roept een vraag op; analisten beoordelen de betekenis. Waarschuwingslijsten helpen mogelijk misleidende waarden herkennen, maar geven geen universeel oordeel. Ook verrijking heeft bronnen die beoordeeld moeten worden. Leg uit wat een record ondersteunt en neem niet iedere verbonden bewering over in een incident.}
\end{frame}

\begin{frame}{Kies de actie voordat u de uitvoer kiest}
\begin{itemize}
\item Onderzoeken: gebruik context om betere vragen te stellen.
\item Detecteren of zoeken: zoek relevante waarnemingen in uw data.
\item Blokkeren: pas beoordeeld beleid toe met aanvaardbare impact.
\end{itemize}
\note{Gebruikers en tools verdragen foutpositieven, oude informatie en onzekerheid niet allemaal op dezelfde manier. Een waarneming kan een zoekactie rechtvaardigen zonder geschikt te zijn voor automatisch blokkeren. De to-IDS-markering is onderdeel van exportbeleid en bewijst geen algemene veiligheid of relevantie. Beoordeel bron, betrouwbaarheid, ouderdom, beperkingen en doelsysteem voor operationeel gebruik. Deze inleiding schrijft geen productiebeleid voor blokkadelijsten voor.}
\end{frame}

\begin{frame}{Kies wie de informatie mag ontvangen}
\begin{itemize}
\item Stel distributie en deelgroepen bewust in.
\item Controleer event, attributen en beoogde ontvangers.
\item Gebruik labels zoals TLP naast toegangscontroles.
\item Controleer publicatie en instellingen vóór uitwisseling.
\end{itemize}
\note{Een deelgroep bepaalt een toegestane doelgroep binnen het ingestelde beleid. TLP beschrijft verwachtingen over omgang met informatie; een TLP-tag verleent of handhaaft op zichzelf geen zichtbaarheid in MISP. Publicatie en een ingesteld synchronisatie-, feed- of exportpad zijn andere beslissingen dan het distributieniveau. Lokale koppelingen en eigen kennis kunnen extra grenzen toevoegen. Begin met een ongepubliceerd event dat alleen voor de eigen organisatie zichtbaar is.}
\end{frame}

\begin{frame}{Kwaliteit groeit door feedback}
\begin{itemize}
\item Leg herkomst, betrouwbaarheid en bruikbare tijdcontext vast.
\item Beoordeel treffers op waarschuwingslijsten en verrijking.
\item Voeg waarnemingen, correcties en foutpositieve meldingen toe.
\item Herbeoordeel oude informatie en het beoogde gebruik.
\end{itemize}
\note{Een sighting registreert een waarneming en verhoogt niet automatisch de betrouwbaarheid van iedere interpretatie. Onderscheid waarnemingen, foutpositieven en verval/afloop. Het waarnemingstijdstip verschilt van het wijzigingstijdstip van een event of attribuut. Ouderdom en vervalmodellen ondersteunen de beoordeling, maar vervangen geen lokaal oordeel. Feedback moet terechtkomen bij de mensen of processen die de informatie onderhouden.}
\end{frame}

\begin{frame}{Een deelgemeenschap heeft werkafspraken nodig}
\begin{itemize}
\item Spreek doel, deelnemers en verwachte bijdragen af.
\item Bepaal hoe informatie wordt behandeld, beoordeeld en verbeterd.
\item Begin met een kleine vertrouwde groep en bruikbare informatie.
\item Verbind andere gemeenschappen als het beleid dat toelaat.
\end{itemize}
\note{Een gemeenschap kan een intern team, sectorgroep, nationaal of CSIRT-netwerk of thematische samenwerking zijn. Instanties kunnen zelfstandig werken of geselecteerde informatie uitwisselen. Een centrale openbare gemeenschap of een wereldwijd gebruikersaantal is niet nodig om dit model uit te leggen. De technische opzet schept op zichzelf geen vertrouwen of rechtsgrond. Bespreek verwachtingen met beheerders en deelnemers.}
\end{frame}

\begin{frame}{Maak van een obstakel een eerste stap}
\small
\begin{tabular}{@{}p{.28\textwidth}p{.68\textwidth}@{}}
\toprule
Obstakel & Bruikbare eerste stap \\
\midrule
Weinig tijd & Deel een kleine, gestructureerde waarneming met de bron. \\
Onzekerheid & Benoem wat bekend, onbekend en voor correctie vatbaar is. \\
Verschillende tools & Spreek begrippen af en gebruik ondersteunde uitwisselingsformaten. \\
Verplichtingen & Beoordeel doel, ontvangers en toepasselijk beleid met de verantwoordelijken. \\
\bottomrule
\end{tabular}
\note{Delen vraagt aandacht voor vertrouwen, tijd, motivatie en organisatorische/juridische verplichtingen, naast technologie. Een tool of label maakt delen niet automatisch conform alle regels. Een eerste bruikbare bijdrage kan context of een correctie zijn in plaats van een nieuwe indicatorenlijst. Automatisering moet herhaald werk beperken en tegelijkertijd beoordeling en betekenis behouden.}
\end{frame}

\begin{frame}{Het project begon met een praktische behoefte}
\begin{itemize}
\item Malwareanalisten wilden dubbel werk voorkomen.
\item Praktische feedback liet het vroege platform groeien.
\item CIRCL en een bredere gemeenschap ontwikkelen MISP verder.
\item Code, kennis en opleiding zijn bruikbare bijdragen.
\end{itemize}
\note{De eerdere inleiding beschrijft de malwarewerkgroep uit 2012 en het vroege platformwerk van Christophe Vandeplas. Behoud die oorsprong zonder er een verkooppraatje van te maken. CIRCL draagt al lang bij aan ontwikkeling en gemeenschappen; bijdragers en beheerders brengen bredere toepassingen voor incidentrespons, fraude en analyse mee. Vermijd ongedateerde gebruikersaantallen of onbevestigde claims over institutionele mandaten en marktaandelen. De README verwijst naar projectinformatie.}
\end{frame}

\begin{frame}{Het MISP-project is meer dan één toepassing}
\IntroDiagram{overview}
\note{Vier elkaar aanvullende onderdelen: open software, gedeelde kennis, standaarden en gemeenschappen. MISP, PyMISP en optionele modules ondersteunen gebruik; taxonomieën, Galaxies en waarschuwingslijsten leveren herbruikbare kennis; formaten, begrippen en objectdefinities helpen bij uitwisseling; mensen leveren werkwijzen, opleiding en feedback. De categorieën overlappen en vormen geen hiërarchie of verplichte dienstenlijst. Het oude losse dashboard wordt niet voorgesteld als het huidige geïntegreerde dashboard.}
\end{frame}

\begin{frame}{De gemeenschap geeft MISP vorm}
\IntroDiagram{community}
\note{Deze bewerkbare bijdrage- en feedbackcyclus vervangt het oude bestuursdiagram. Gebruikers en beheerders brengen behoeften in; bijdragers stellen code, kennis, documentatie en opleiding voor; projectbeheerders beoordelen, testen en besluiten over releases; verbeteringen keren terug naar de praktijk. Mensen kunnen meerdere rollen hebben en ieder repository heeft een eigen bijdrageproces. Dit is een conceptuele uitleg, geen formeel juridisch of organisatorisch gezagsschema. Volg de richtlijnen van het betreffende repository.}
\end{frame}

\begin{frame}{Eerste oefening: leg een bruikbaar record uit}
\begin{enumerate}
\item Maak een privé, ongepubliceerd event voor de fictieve casus.
\item Voeg domein en IP toe met bron, tijd en onzekerheid.
\item Leg een mogelijke correlatie en passende vervolgstap uit.
\item Beoordeel ontvangers en omgang vóór publicatie.
\end{enumerate}
\note{De oefening van vijf minuten kan op papier en vereist geen MISP-account. Gebruik op een toegestane trainingsinstantie alleen wegwerpdata en houd het event ongepubliceerd en beperkt tot de eigen organisatie. Vraag wat Team B aan het record heeft, wat onbewezen blijft en hoe een correctie of waarneming terugkomt. Publiceer of blokkeer de documentatiewaarden niet als onderdeel van de inleiding.}
\end{frame}

\begin{frame}{Ga verder met de taak die u nodig hebt}
\begin{itemize}
\item Dagelijkse analyse: MISP-gebruik en tags.
\item Meer context: objecten, Galaxies en informatiekwaliteit.
\item Integratie: de workshops REST-API en PyMISP.
\item Gemeenschappen beheren: administratie en deelpraktijken.
\end{itemize}
\note{De README bevat directe links naar de vervolgcursussen en het trainingsoverzicht. Deze inleiding levert een gemeenschappelijke woordenschat en een besluitvormingsproces; configuratie en API-verzoeken horen in de praktische modules. De diagrammen delen TikZ-geometrie en hebben vertaalde labels en losse vectorexports, zodat trainers ze zonder opnieuw tekenen kunnen aanpassen.}
\end{frame}

\begin{frame}{Deel bruikbare context en blijf die verbeteren}
\begin{itemize}
\item Structureer de waarneming en behoud de betekenis.
\item Kies doelgroep en actie bewust.
\item Gebruik feedback om kwaliteit en begrip te verbeteren.
\end{itemize}
\note{Laat deelnemers de casus in eigen woorden uitleggen: wat delen zij, met wie, voor welk gebruik en hoe kan het worden gecorrigeerd? Het doel is een bruikbare deelpraktijk met ondersteuning van MISP. Het open project combineert software, kennis, standaarden en gemeenschappen. Deelname kan beginnen met een kleine bijdrage die een andere analist helpt beter te beslissen.}
\end{frame}
